\section{解消族を通じた一つの Kähler corner 模型}
\label{tc:section}

第\ref{sc:phase}節の半安定チャートと第\ref{exc:section}節の曲面模型を、
一つの明示的な Kähler form の中で実現する。
Hamiltonian torus の moment と不変な複素関数を組み合わせる手法自体は、
Gross \cite[\S\S1--2]{gross-toric}、Auroux \cite[\S3.3.1]{auroux-blowup}、
Casta\~no-Bernard--Matessi \cite[\S\S3.4,5.1]{cbm-conifold} の既知の構成に属する。
ここで必要なのは、今回の ordered resolution、二つの exceptional areas、
中心成分上の exceptional sphere、非零ファイバーの moments を同じ式で追うことである。
所定の全 boundary data を保つ任意の global Kähler form との比較までは主張しない。

\subsection{二つの射影直線への graph closure}

$V=\{xyz=tw\}\subset\CC^5$ とする。
有理写像 $[x:t]$, $[xy:t]$ の graph closure を $\widetilde V$ と書く。
二つの homogeneous coordinates を $[a:b]$, $[c:d]$ とすると、その方程式は
\begin{equation}\label{tc:incidence}
 xb=ta,\qquad bc=ay\,d,\qquad cz=dw.
\end{equation}

\begin{proposition}\label{tc:resolution}
$\widetilde V$ は滑らかな複素四次元多様体であり、$V$ への proper birational map を持つ。
$t:\widetilde V\to\CC$ は第\ref{sc:phase}節の ordered semistable resolution と同じである。
具体的には次の三つの affine charts で覆われる。
\[
\begin{array}{c|c|c}
\text{chart}&\text{coordinates}&\text{relations}\\ \hline
 b\ne0 &(t,X,y,z)&x=tX,\quad w=Xyz\\
 a,d\ne0 &(x,T,Y,z)&t=xT,\quad y=TY,\quad w=Yz\\
 a,c\ne0 &(x,y,U,w)&t=xyU,\quad z=Uw.
\end{array}
\]
原点上の exceptional fiber は二本の射影直線 $\ell_1\cup\ell_2$ であり、
二つの projection line bundles のそれぞれの次数は
\[
\begin{array}{c|cc}
 &\ell_1&\ell_2\\ \hline
 \mathrm{pr}_1^*\mathcal O(1)&1&0\\
 \mathrm{pr}_2^*\mathcal O(1)&0&1.
\end{array}
\]
\end{proposition}

\begin{proof}
$b=1$ とすると、第2式から $d\ne0$ である。
$d=1,a=X$ と置けば $c=Xy$ となり、最初の chart を得る。
$a=d=1,b=T,c=Y$ と置くと第二の chart を得る。
$a=c=1,d=U$ と置くと $b=yU$ となり、第三の chart を得る。
$b=0$ なら $a\ne0$ なので、残る二つの charts がその部分を覆う。
各 chart は $\CC^4$ であり、もとの有理写像の graph と稠密開集合で一致する。
従って式\eqref{tc:incidence}は余分な成分を持たず、graph closure を定める。
Projection は projective なので proper である。

第二の chart で、原点上の fiber は
$\{x=z=Y=0\}$ と $\{x=z=T=0\}$ の和である。
前者では $[a:b]=[1:T]$ のみ、後者では $[c:d]=[Y:1]$ のみが動く。
両者を隣の charts で閉じれば主張する二本の $\PP^1$ と次数表を得る。
\end{proof}

$\int_{\PP^1}\omega_{\mathrm{FS}}=1$ とし、$e_1,e_2>0$ に対し
\begin{equation}\label{tc:kahler}
 \widetilde\Omega
 =\omega_{\CC^5}
   +e_1\mathrm{pr}_1^*\omega_{\mathrm{FS}}
   +e_2\mathrm{pr}_2^*\omega_{\mathrm{FS}}
 \quad\text{on }\widetilde V
\end{equation}
と置く。これは product Kähler form の埋め込みによる制限なので Kähler である。
また $\int_{\ell_i}\widetilde\Omega=e_i$ である。
Ordered indices $i<j<k$ に対して $e_1=(j-i)\eps$, $e_2=(k-j)\eps$ とすれば、
既述の polarization $aH+\eps A$ が原点上の二本の曲線に与える areas と一致する。
六次の $k<5<6$ では $e_1=(5-k)\eps$, $e_2=\eps$ である。

\subsection{非零ファイバー上の proper な torus fibration}

$t=\tau\ne0$ を固定すると、$\widetilde V_\tau\simeq\CC^3_{x,y,z}$ であり
$w=xyz/\tau$ となる。
角は周期1、$\rho_u=\pi|u|^2$ と規格化する。
次を置く。
\begin{align}
 A_1&=e_1\frac{|x|^2}{|x|^2+|\tau|^2},&
 A_2&=e_2\frac{|xy|^2}{|xy|^2+|\tau|^2},\nonumber\\
 M_1&=\rho_x-\rho_y+A_1,&
 M_2&=\rho_y-\rho_z+A_2.\label{tc:moments}
\end{align}
$c>0$ に対し $Q=|w-c|^2$ とする。

\begin{theorem}\label{tc:fibration}
$F_\tau=(M_1,M_2,Q)$ は $Q>0$ 上の proper な Lagrangian $T^3$-fibration である。
ここで proper は、target $\RR^2\times\RR_{>0}$ の compact subsets の逆像が compact という意味である。
Discriminant は $Q=c^2$ にある三本の rays
\begin{equation}\label{tc:discriminant}
 \{(m,0):m\ge0\},\quad
 \{(-m,m):m\ge0\},\quad
 \{(0,-m):m\ge0\}
\end{equation}
である。各 open ray では transverse focus--focus 型であり、
singular fiber は $(\text{pinched }T^2)\times S^1$ である。
原点の vertex fiber は $T^3$ の一つの $T^2$ を一点へ潰した空間である。
\end{theorem}

\begin{proof}
座標の rotations による $T^3$-action は $w$ も回すので、
その三つの moments は
\[
 p_x=\rho_x+\rho_w+A_1+A_2,\quad
 p_y=\rho_y+\rho_w+A_2,\quad
 p_z=\rho_z+\rho_w.
\]
積 $xyz$ を固定する $T^2$ の generators を $(1,-1,0),(0,1,-1)$ とすれば、
対応する moments は $M_1=p_x-p_y,M_2=p_y-p_z$ である。
$w$ はこの $T^2$ に不変なので、三つの関数は Poisson commute する。

まず $(M_1,M_2,w)$ が orbit space を $\RR^2\times\CC$ と同定することを示す。
$w\ne0$ では $s_i=\log|z_i|^2$ とし、$\sum s_i$ を固定する。
Moments は、strictly convex な関数
\[
 G(s)=\pi\sum_i e^{s_i}
 +\frac{\pi e^{s_x+s_y+s_z}}{|\tau|^2}
 +e_1\log(e^{s_x}+|\tau|^2)
 +e_2\log(e^{s_x+s_y}+|\tau|^2)
\]
の gradient の差である。
指定した $M_1,M_2$ に対応する線形関数を引き、平面 $\sum s_i=\mathrm{const}$ に制限する。
この平面上で $|s|\to\infty$ なら少なくとも一つの $s_i\to+\infty$ なので、
最初の指数関数項が線形関数に勝ち、関数は coercive である。
従って unique minimizer があり、指定した moments と $|w|$ に対する radii は一意に存在する。
位相の積を指定すれば、残りは一つの $T^2$-orbit である。

$w=0$ では、三つの coordinate planes の images はそれぞれ
\[
\begin{array}{c|c}
z=0&M_2\ge0,\ M_1+M_2\ge0\\
y=0&M_1\ge0,\ M_2\le0\\
x=0&M_1\le0,\ M_1+M_2\le0.
\end{array}
\]
これらは $\RR^2$ を覆い、共通境界は式\eqref{tc:discriminant}の rays である。
各 plane の toric moment map は quadrant 上に一対一かつ全射である。
例えば $z=0$ では $p_x=M_1+M_2,p_y=M_2$ を固定し、
$G(s_x,s_y)-p_xs_x-p_ys_y$ を最小化すれば、quadrant の内部について同じ strict convexity と coercivity が使える。
Boundary では対応する coordinate を零とし、残る一変数の strictly increasing moment を解けばよい。
従って $w=0$ でも一つの orbit が定まる。

Moments と $w$ が有界なら、$0\le A_i\le e_i$ より三つの $\rho_i$ の差は有界である。
一つの $\rho_i$ が無限大になるなら全てが無限大になるが、
\[
 \rho_x\rho_y\rho_z=\pi^3|\tau|^2|w|^2
\]
に反する。従って $(M_1,M_2,w)$ は proper な orbit-space map であり、
上の全単射は homeomorphism を与える。
同じ評価により $F_\tau$ も proper である。

$T^2$-action は少なくとも二つの coordinates が非零なら free である。
その locus では $(M_1,M_2,w)$ は smooth な quotient coordinates となる。
$w\ne0$ の場合は上の strict convexity、coordinate plane の smooth points では
残る coordinate が $w$ の transverse complex coordinate になることから従う。
$Q>0$ では $|w-c|^2$ は reduced surface 上で regular なので、
正則 fiber は circle 上の principal $T^2$-bundle である。
この bundle は自明であり、fiber は compact connected $T^3$ となる。
Commutativity と rank により Lagrangian である。

非零 $x$-axis の stabilizer は $(0,1,-1)$ で、normal weights は $(1,-1)$ である。
$Q$ の normal quadratic は
\[
 -2\operatorname{Re}(\bar c\,x yz/\tau)
\]
であり、その係数は非零である。
Invariant Kähler metric によって normal lines を規格化すれば、
stabilizer moment とこの quadratic は focus--focus Cartan を与える。
他の二軸も同様である。
Axis の moment images は直接式\eqref{tc:moments}から式\eqref{tc:discriminant}となる。
Reduced circle が $w=0$ を通る一箇所で primitive $S^1\subset T^2$ が潰れるため、
open ray の fiber は主張した積空間となる。
Vertex では $w=0$ の orbit が一点であり、その一箇所で全 $T^2$ が潰れる。
Circle をその点で切った interval 上で torus bundle を自明化すれば、
最後の topological description を得る。
\end{proof}

この定理は、標準 $\CC^3$ の positive-vertex topology を、
二つの exceptional periods を持つ $\widetilde\Omega|_{\widetilde V_\tau}$ の上で実現している。
Toric terms を無視して標準 Euclidean moments に置き換える必要はない。
$Q=0$ の elliptic locus は本節の fibration domain から外している。
また bounded base window で切り取った場合の outer boundary は、この局所定理では埋めていない。

\subsection{中心成分上では exceptional bridge に厳密に一致する}
\label{tc:surface}

六次の $k<5<6$ を考え、第二の chart $t=xT,y=TY,w=Yz$ を使う。
$D_5=\{T=0\}$ 上では、式\eqref{tc:kahler}は
\begin{equation}\label{tc:central-form}
 \widetilde\Omega|_{D_5}
 =\omega_x+\omega_z+\omega_{Yz}+e_2\omega_{\mathrm{FS}}[Y:1]
 =\omega_x+\omega_{e_2}
\end{equation}
という product form になる。第一の projective factor は $[1:0]$ で一定である。
従って $S_{k5}=\{x=T=0\}$ 上には、全 exceptional sphere を含む
第\ref{exc:section}節の form がそのまま現れる。
Charts を $Y=\infty$ まで継ぐと、これは $\CC_x\times\operatorname{Bl}_0\CC^2$ 上の同じ式である。

Torus action と moments は全 $\widetilde V$ に延びる。
$D_5$ 上での制限は
\begin{equation}\label{tc:central-moments}
 M_1=e_1+\rho_x,\qquad
 M_2=e_2-J,\qquad
 Q=|Yz-c|^2,
 \qquad J=\rho_z+\frac{e_2}{1+|Y|^2}.
\end{equation}
第一の circle は $(x,T,Y,z)$ に weights $(1,-1,0,0)$ で作用するので、
$D_5$ 内では $S_{k5}$ を点ごとに固定する所定の normal circle である。
第二の circle は $(Y,z)$ に weights $(1,-1)$ で作用し、$J$-circle の逆向きである。
従って $C_{56}$ の端 $Y=z=0$ は $(M_1,M_2)=(e_1,0)$ にあり、
exceptional sphere はそこから strict triple boundary の $(e_1,e_2)$ へ写る。
この線分の内部は reduced-surface fibration の regular fibers と交わるもので、
非零ファイバーの discriminant edge という意味ではない。

\begin{corollary}\label{tc:common-model}
選んだ form \eqref{tc:kahler}に対して、次の三つは一つの Kähler degeneration の中で両立する。
\begin{enumerate}
\item 非零ファイバー上の compact positive-vertex torus fibers。
\item 中心成分上の、exceptional sphere 全体を通る Auroux blow-up model。
\item $C_{56}$ の一端における所定の divisor-normal circle と、
      第\ref{exc:section}節の inverse-radial focus family。
\end{enumerate}
$x$ が零から離れた compact annulus と、$D=\{U=0\}$ を除いた第二の chart の
compact subsets 上では、非零ファイバーの forms と maps は
式\eqref{tc:central-form}--\eqref{tc:central-moments}へ $C^\infty$ に収束する。
これは $D$ を含む全 exceptional sphere への一様収束を主張するものではない。
\end{corollary}

\begin{proof}
第一項は定理\ref{tc:fibration}、第二・第三項は二つの制限式から従う。
Annulus 上では $T=\tau/x$, $y=\tau Y/x$ を代入する。
これらと全 derivatives は $\tau\to0$ で零へ収束し、projective coordinates も滑らかに収束する。
したがって compact subsets 上の $C^\infty$ 収束を得る。
Whole compact fiber を扱う場合は、その全体を含む compact subset を先に選ぶ。
特に compact focus fiber は $J=e_2>0$ なので $D$ を避け、この選択ができる。
他方、第三の chart で $y=\tau/(xU)$ と解く表示は $U=0$ で一様にならない。
\end{proof}

\subsection{三つの成分の residue を保つ logarithmic primitive}
\label{tc:log-primitive}

三つの charts を順に $A,B,C$ とし、$\lambda_A^0,\lambda_B^0,\lambda_C^0$ を
式\eqref{tc:kahler}の各 affine chart 上の標準 primitive とする。
すなわち Euclidean terms には $\rho_u d\theta_u$ を、
projective affine coordinate $v$ には
$e|v|^2(1+|v|^2)^{-1}d\theta_v$ を用いる。
これらは各 affine chart の coordinate axes まで滑らかである。
Charts $A,B$ の第一の projective coordinates は $X,T=1/X$ であり、
charts $B,C$ の第二のものは $Y,U=1/Y$ である。
従って overlap 上で厳密に
\begin{equation}\label{tc:primitive-changes}
 \lambda_B^0=\lambda_A^0-e_1d\theta_X,
 \qquad
 \lambda_B^0=\lambda_C^0-e_2d\theta_U
\end{equation}
となる。

\begin{proposition}[同時に貼れる正の residues]
\label{tc:primitive-proposition}
Ordered components を $D_i,D_j,D_k$ とし、任意の $w_i>0$ に対し
\[
 w_j=w_i+e_1,\qquad w_k=w_j+e_2
\]
と置く。次の三つの表示は、$\widetilde V\setminus\{t=0\}$ 上の一つの smooth primitive
$\lambda$ を定め、$d\lambda=\widetilde\Omega$ を満たす。
\begin{equation}\label{tc:primitive-global}
\begin{aligned}
 A:\quad&\lambda_A^0-w_i d\theta_t,\\
 B:\quad&\lambda_B^0-w_i d\theta_x-w_jd\theta_T,\\
 C:\quad&\lambda_C^0-w_i d\theta_x-w_jd\theta_y-w_kd\theta_U.
\end{aligned}
\end{equation}
中心成分への limiting angular residues はそれぞれ $-w_i,-w_j,-w_k$ である。
非零ファイバーに制限すると、$T^2$ generators $u,v$ について
$\lambda(u)=M_1$, $\lambda(v)=M_2$ となる。
\end{proposition}

\begin{proof}
$A,B$ の重なりでは $x=tX,T=X^{-1}$ なので、
式\eqref{tc:primitive-changes}を使うと $B$ の表示は
\[
 \lambda_A^0-w_i d\theta_t
       +(w_j-w_i-e_1)d\theta_X
\]
となり、最後の係数は零である。
$B,C$ の重なりでは $T=yU$ なので、$B$ の表示は
\[
 \lambda_C^0-w_i d\theta_x-w_jd\theta_y
       -(w_j+e_2)d\theta_U
\]
となり、これは $C$ の表示である。
従って三つは厳密に貼り合う。
各 angular term の微分は complement 上で零であり、各 $\lambda^0$ は
$\widetilde\Omega$ の primitive なので、微分の主張も従う。
各 coordinate divisor に沿う $\lambda^0$ の meridian integral は零へ収束するため、
residues は指定した値となる。

$t=\tau\ne0$ の fiber は全て chart $A$ に含まれる。
そこでは $d\theta_t=0$ なので、$\lambda$ は制限 form の toric primitive そのものである。
その contractions は $p_x-p_y,p_y-p_z$ であり、最後の主張を得る。
\end{proof}

この式は、component-normal directions を二つだけ残すと見失う residue を具体化する。
$S_{ij}$ 上の第二の exceptional sphere では normal degrees は $(0,-1)$ であり、
三つ目の divisor は一点で交わる。
第\ref{ct:degree-residue}項の収支は
\[
 w_k-w_j=e_2
\]
となり、上の primitive はこの収支を正の weights で満たしている。
二つの normal directions だけで全 sphere に smooth な $\beta$ を要求する場合の
$e_2+w_j=0$ という不可能な条件とは、入力が異なる。
ただし logarithmic primitive が貼れたことだけから、
その Liouville flow が $F_\tau$ の全 compact fibers を保存することは従わない。
所定の flow に適合する foliation と端の比較は、依然として moving-family の問題である。

\begin{remark}\label{tc:nonprojectable}
実際、この slice primitive の Liouville field $Z$ は $F_\tau$ に関して一般には projectable でない。
Dense torus では $\lambda=\sum p_i d\theta_i$ なので $Zp_i=p_i$ である。
$M_1=M_2=0$ では $p=P(1,1,1)$、$P>0$ であり、上の strictly convex potential の Hessian を
$\mathcal H$ とすると
\[
 Z\log|w|=\frac P2\,\boldsymbol1^{\mathsf T}\mathcal H^{-1}\boldsymbol1>0.
\]
$Z$ は toric angles を保つ。
同じ正則 fiber $M_1=M_2=0,Q=q^2$、$0<q<c$ にある二つの点
$w=c+q$ と $w=c-q$ では、$ZQ=2w(w-c)Z\log|w|$ の符号が反対となる。
従って $ZQ$ はその fiber 上で一定でない。
これは primitive の整合と Liouville に適合する fiber decomposition の整合を区別する具体例である。
\end{remark}

\subsection{この構成で残る比較問題}

式\eqref{tc:kahler}は実際の局所解消空間にある明示的な form であり、
単に abstract な blow-up の位相模型を与えたものではない。
ただし exceptional curves の areas が所定の polarization と一致することから、
任意に先に指定した global form と全 collar foliations を同時に保つ同定が直ちに従うわけではない。
その結論には、local Kähler representative の貼り合わせ、または必要な strata・normal first jets・
whole fibers を保持する relative symplectic comparison が必要である。
また中心成分での厳密な一致と、$\tau\ne0$ における既存 pants block への厳密な接着も区別する。
後者では family の return action と end transport をさらに合わせる必要がある。

この common model の central framing における focus winding は零である。
第\ref{cw:section}節の結果により、smooth equivariant normal collars のまま
全24端をこの型に揃えると、determinant degree 36 の収支を満たせない。
従って本節は一つの整合した corner-to-curve \emph{局所模型}を完成させるが、
全端の relative gluing を完成させたものではない。
